\documentclass[11pt]{article}
\usepackage[margin=1.05in]{geometry}
\usepackage{amsmath,amssymb,amsthm,mathrsfs}
\usepackage{hyperref}

\newtheorem{thm}{Theorem}[section]
\newtheorem{lem}[thm]{Lemma}
\newtheorem{prop}[thm]{Proposition}
\newtheorem{cor}[thm]{Corollary}
\theoremstyle{definition}
\newtheorem{defn}[thm]{Definition}
\newtheorem{rem}[thm]{Remark}
\newtheorem{ex}[thm]{Example}

\newcommand{\Z}{\mathbb{Z}}
\newcommand{\Saff}{\widetilde{S}_{n}}
\newcommand{\cont}{\operatorname{content}}
\newcommand{\et}{\tilde e_1}
\newcommand{\XX}{\mathcal{X}}
\newcommand{\EE}{\mathcal{E}}
\newcommand{\BB}{\mathcal{B}}

\title{Every Morse--Schilling crystal edge is an exchange move,\\
and the exchange move is strictly more general}
\author{Clio}
\date{22 September 2026}

\begin{document}
\maketitle

\begin{center}\small
\begin{tabular}{ll}
\textbf{Author:} Clio &
\textbf{Date:} 22 September 2026 \\
\textbf{Recipient:} Robin Langer, cc Rick (grandpa-rick) &
\textbf{Commit:} \texttt{clio-vega/proofs@PENDING} \\
\end{tabular}\\[2pt]
\textbf{Title:} Every Morse--Schilling crystal edge is an exchange move, and the
exchange move is strictly more general.\\[1pt]
\emph{The commit named above contains the mathematics of this note; nothing below this
cover block has changed since.}
\end{center}

\begin{abstract}
Morse and Schilling \cite{MS} equip the set of two-factor affine factorisations of
$w\in\Saff$ with a $U_q(\mathfrak{sl}_2)$-crystal, whose raising operator $\et$ is defined
after fixing a parameter $x\in[n]$ with $x\notin\cont(u)\cup\cont(v)$ and pairing the two
contents with respect to $x$. In \cite{Clio2} we gave an exchange move on the same objects
which uses no such parameter. We settle the relation between the two. Let $\BB(S,T)$ be the
set of \emph{all} $\et$-images of $u_Su_T$, taken over \emph{every} admissible $x$ (and, when
no $x$ is admissible, over every instance of the reduction of \cite[\S\texttt{two factor case}]{MS}), and let
$\EE(S,T)$ be the set of exchange moves of \cite[Thm.~4.1]{Clio2}. We prove
$\BB(S,T)\subseteq\EE(S,T)$ for every length-additive $(S,T)$ (Theorem~\ref{thm:contain}),
and we show the containment is strict, the smallest witness being $n=5$
(Theorem~\ref{thm:strict}). We then characterise the difference exactly
(Theorem~\ref{thm:charac}): writing the factorisation as a bracket word and letting $Q$ be
its slot-wise low-point profile, membership in $\EE$ is \emph{local} minimality of $Q$ --- the
last point of the plateau that opens a run of $S$ --- whereas membership in $\BB$ is
\emph{global} minimality of $Q$ over a full period beginning at an admissible cut. The
$x$-pairing is therefore not a removable decoration: it is exactly the device that promotes
local minimality to global minimality, and it buys canonicity (one well-defined operator,
invertible, satisfying the Stembridge axioms) at the price of coverage. The intended reading
of Theorem~4.1 as ``the $x$-free shadow of the whole family of crystal structures'' is
\emph{false} in the direction one would guess: the union over all $x$ is strictly smaller,
not equal.
\end{abstract}

\section{Introduction}

\subsection{The question, and two of my own sentences in tension}

Let $n=k+1\ge2$, let $\Saff$ be the affine symmetric group, and for a proper subset
$S\subsetneq\Z/n$ let $u_S$ be the unique cyclically decreasing element with support $S$
\cite[Lem.~2.2]{Clio2}. Call a pair $(S,T)$ of proper subsets \emph{additive} if
$\ell(u_Su_T)=|S|+|T|$; equivalently $u_Su_T$ is a two-factor affine factorisation in the
sense of Lam \cite{Lam2006}, with $\cont(u)=S$ and $\cont(v)=T$ in the notation of \cite{MS}.

Morse and Schilling \cite[Def.~\texttt{crystal operators}]{MS} define $\et$ on such a pair after fixing
$x\in[n]$ with $u,v\in S_{\hat x}$, i.e.\ $x\notin S\cup T$; write
\[
  \XX(S,T)\;=\;\Z/n\setminus(S\cup T)
\]
for the set of admissible parameters. Theorem 4.1 of \cite{Clio2} gives a move on the same
objects and fixes no parameter. The question this note answers was posed as: is the exchange
move the $x$-free shadow of the whole family $\{\et^{(x)}\}_{x\in\XX}$, i.e.\ is
\[
  \bigcup_{x\in\XX(S,T)}\bigl\{\et^{(x)}(uv)\bigr\}\setminus\{\mathbf 0\}
  \;\overset{?}{=}\;\bigl\{\text{Theorem 4.1 moves on }uv\bigr\}.
\tag{Q}
\]

Two sentences in my own notes of 21 September were in tension, and the first business of this
note is to write their quantifiers out.

\begin{itemize}
\item[(a)] ``$\{\et\text{ edges}\}\subsetneq\{\text{my moves}\}$'', recorded after an
exhaustive comparison at a \emph{single} $x$.
\item[(b)] ``$|L_1|-|R_1|=|S|-|T|$, so my hypothesis $|S|>|T|$ is strictly stronger than
Morse--Schilling's firing condition $L_1\ne\emptyset$; hence $\et$ is never $\mathbf 0$ in my
regime'' --- from which I had concluded that $\et$ can fire where Theorem 4.1 offers nothing,
contradicting (a) read as an unrestricted containment.
\end{itemize}

Both statements are true and they do not conflict. Statement (b) is not an empirical fact but a
one-line consequence of the definition: the $uv$-pairing is a partial matching between $S$ and
$T$, so if $\pi$ denotes the number of pairs then $|L_1|=|S|-\pi$ and $|R_1|=|T|-\pi$, whence
$|L_1|-|R_1|=|S|-|T|$ identically, for every $x$. The inference drawn from it was the error,
and the broken assumption was about \emph{my own} theorem: the hypothesis $|S|>|T|$ of
\cite[Thm.~4.1]{Clio2} is used \emph{only} in the applicability lemma
\cite[Lem.~4.2]{Clio2}, which asserts that a usable run exists. The move itself
\cite[Lem.~4.3]{Clio2} requires only that a usable run \emph{be given}, together with
$|T|\le n-2$. So the object to compare with $\BB$ is not ``the moves Theorem 4.1 guarantees''
but ``the moves Theorem 4.1 constructs'':

\begin{defn}\label{def:E}
Let $(S,T)$ be additive. A run $[m,M]$ of $S$ is \emph{usable} if $m\notin T$ and $j+1\notin T$
for some $j\in[m,M]$. For a usable run put $e_{[m,M]}=$ the first $j\in\{m,m+1,\dots,M\}$ with
$j+1\notin T$. Set
\[
  \EE(S,T)=\bigl\{\,e_{[m,M]} \;:\; [m,M]\text{ a usable run of }S \,\bigr\}\subseteq S ,
\]
and write $\EE^{\mathrm{mv}}(S,T)=\bigl\{(S\setminus\{e\},\,T\cup\{m\}):e=e_{[m,M]}\bigr\}$
for the corresponding moves.
\end{defn}

With this reading (a) is unrestricted after all, and (b) costs nothing: there is no regime in
which $\et$ fires and Theorem 4.1 is silent. That is the content of Theorem~\ref{thm:contain}.
Once the comparison is set up this way, (Q) has a clean answer, and it is not the answer the
question anticipated.

\subsection{The answer}

Write
\[
  \BB(S,T)=\bigl\{\,b_x \;:\; x\in\XX(S,T),\ \et^{(x)}(uv)\neq\mathbf 0\,\bigr\}\subseteq S,
  \qquad b_x=\min\nolimits_{\prec_x} L_1^{(x)}(uv),
\]
so that $\et^{(x)}(uv)$ removes $b_x$ from $S$ and adds to $T$ the bottom of the run of $S$
containing $b_x$. Then:

\begin{itemize}
\item $\BB(S,T)\subseteq\EE(S,T)$ always, and the two moves agree letter for letter
(Theorem~\ref{thm:contain}). Every edge of every Morse--Schilling crystal on $uv$, at every
admissible $x$, is an exchange move.
\item The containment is strict. The smallest witnesses occur at $n=5$; at $n=5,6,7$ there are
respectively $15$, $132$ and $854$ exchange moves lying outside every Morse--Schilling crystal
(Theorem~\ref{thm:strict} and \S\ref{sec:verify}).
\item The difference is exactly the gap between local and global minimality of a height
profile (Theorem~\ref{thm:charac}).
\end{itemize}

So the sentence one would have liked --- ``Theorem 4.1 is the $x$-free shadow of the whole
family of Morse--Schilling crystals, and dropping the pairing is taking the union over all of
them'' --- is false. The true sentence is one-directional and, I think, more informative:

\begin{quote}
\emph{The exchange move is strictly more general than the union of the entire family of
Morse--Schilling crystal structures on $uv$. The $x$-pairing is what turns a local minimum
into a global one; it buys a canonical, invertible operator and it costs coverage.}
\end{quote}

\begin{rem}
Neither side of this comparison was built with the other in view, and neither has a free
parameter I could tune. Morse--Schilling's $x$ is \emph{their} parameter and I quantify over
all of it; the exchange move has none. The implementation of \cite[Def.~\texttt{crystal operators}]{MS} used
below reproduces all three of the worked examples printed in \cite{MS} exactly, for both the
pairing and $\et$ (\S\ref{sec:verify}).
\end{rem}

\section{The two constructions}\label{sec:setup}

We use the block model of \cite[\S2]{Clio2}\footnote{We cite \cite{MS} by the internal names its arXiv e-print source gives to
its environments and subsections (\texttt{ex:pairing} at source line 1288,
\texttt{definition.crystal operators} at 1330, \texttt{theorem.crystal} at 1429,
\texttt{subsection.two factor case} at 1752, \texttt{lemma.existence of i} at 1781,
\texttt{theorem.two factor} at 1822), rather than by displayed number. The displayed numbering
differs between versions --- my own notes of 21 September record the two crystal theorems as
``Theorems 4.7 and 5.5'', where the e-print I hold would print 3.5 and 3.14 --- and a number
that is right for one version is a silent error in the other.}: for $S\subsetneq\Z/n$, identified with its
$n$-periodic lift, $u_S(j)=j-1$ if $j-1\in S$ and $u_S(j)=\min\{c\ge j:c\notin S\}$ otherwise;
$\ell(u_S)=|S|$. A \emph{run} of $S$ is a maximal interval $[m,M]\subseteq\Z$ contained in $S$.
Products act by $(ab)(j)=a(b(j))$, so $u_Su_T$ means ``$u_T$ first''; this matches the factor
order of \cite{MS}, in which $\et$ moves a letter from the left factor to the right.

\subsection{The Morse--Schilling pairing}

Fix $x\in\XX(S,T)$. Morse--Schilling order $\Z/n\setminus\{x\}$ linearly by
\begin{equation}\label{eq:order}
  x-1\succ x-2\succ\dots\succ0\succ n-1\succ\dots\succ x+1 ,
\end{equation}
and pair the $\prec_x$-largest $b\in S$ with the $\prec_x$-smallest $a\in T$ with $a\succ_x b$,
ignoring already-paired letters, proceeding downwards through $S$. $L_1(uv)$ and $R_1(uv)$ are
the unpaired elements of $S$ and of $T$. Then \cite[Def.~\texttt{crystal operators}(i)]{MS}
\[
  \et(uv)=\tilde u\tilde v,\quad
  \cont(\tilde u)=S\setminus\{b\},\quad \cont(\tilde v)=T\cup\{b-t\},
\]
for $b=\min_{\prec_x}L_1(uv)$ and $t=\min\{i\ge0: b-i-1\notin S\}$; and $\et(uv)=\mathbf 0$ if
$L_1(uv)=\emptyset$. Note that $b-t$ is precisely the bottom of the run of $S$ containing $b$,
which is the same letter Definition~\ref{def:E} adds to $T$. The two constructions therefore
differ only in \emph{which letter is removed from $S$}, and the whole question is which
elements of $S$ each of them selects.

\subsection{The word and its profile}

\begin{defn}\label{def:word}
Fix $x\in\XX(S,T)$ and let the \emph{window} be the interval of integers
$(x,x+n]=\{x+1,\dots,x+n\}$, a lift of $\Z/n$ with $x$ itself omitted. Read the slots of the
window in increasing order; at slot $z$ emit the letter $\texttt{)}$ if $z\in T$, and then the
letter $\texttt{(}$ if $z\in S$. The resulting word is $W_x$. Define the \emph{profile}
$P_x,Q_x$ on the window by
\[
  P_x(x)=0,\qquad Q_x(z)=P_x(z-1)-[z\in T],\qquad P_x(z)=Q_x(z)+[z\in S] ,
\]
so $Q_x(z)$ is the height of $W_x$ at the low point of slot $z$ and $P_x(z)$ its height at the
end of slot $z$.
\end{defn}

Since $P_x$ and $Q_x$ obey a recursion independent of $x$ and differ only in their
initialisation, for two cuts $x,x'$ and any $z$ in both windows one has
$Q_x(z)-Q_{x'}(z)=\text{const}$; all statements below compare values of $Q$ within one window
and are therefore independent of the normalisation. Writing $d=|S|-|T|$ we have
$Q(z+n)=Q(z)+d$.

\begin{lem}[The pairing is bracket matching]\label{lem:bracket}
For every $x\in\XX(S,T)$, the $uv$-pairing with respect to $x$ coincides with the standard
bracket matching of $W_x$: $b\in S$ is paired with $a\in T$ if and only if the letter
$\texttt{(}$ at slot $b$ is closed by the letter $\texttt{)}$ at slot $a$.
\end{lem}

\begin{proof}
Induction on $|S|$, the case $S=\emptyset$ being vacuous. Let $b_{\max}$ be the $\prec_x$-largest
element of $S$, so the $\texttt{(}$ at slot $b_{\max}$ is the last $\texttt{(}$ of $W_x$. If
some $\texttt{)}$ follows it, let $a$ be the first such; there is no letter strictly between
them, since any intervening letter would be a $\texttt{(}$, contradicting maximality (within a
slot, $\texttt{)}$ precedes $\texttt{(}$, so the $\texttt{)}$ at slot $a$ is the letter
immediately after the $\texttt{(}$ at slot $b_{\max}$ whenever $a$ is the first
$\texttt{)}$ after it). Morse--Schilling process $b_{\max}$ first and pair it with the
$\prec_x$-smallest free $a\succ_x b_{\max}$ in $T$, which is this $a$; the bracket matching
closes an adjacent $\texttt{()}$, which is this same pair. If no $\texttt{)}$ follows
$b_{\max}$ then Morse--Schilling leave $b_{\max}$ unpaired and the stack never pops it, so it is
unmatched in the bracket matching too. In either case delete the matched pair (or the terminal
$\texttt{(}$) from $W_x$: deleting an adjacent $\texttt{()}$ changes neither the bracket
matching of the remaining letters nor the subsequent steps of the Morse--Schilling procedure,
which continues with the next $b$ and the remaining free letters of $T$. Apply the inductive
hypothesis.
\end{proof}

\begin{lem}[Unmatched letters read off the profile]\label{lem:unmatched}
Let $x\in\XX(S,T)$ and $z_0\in S$ in the window. The letter $\texttt{(}$ at slot $z_0$ is
unmatched in $W_x$ if and only if $Q_x(z)>Q_x(z_0)$ for every $z$ with $z_0<z\le x+n$.
\end{lem}

\begin{proof}
The $\texttt{(}$ at $z_0$ is matched precisely when the height of $W_x$ returns to its value
$Q_x(z_0)$ immediately before that letter at some later point of the word. The heights attained
after that letter are $P_x(z_0)=Q_x(z_0)+1$ together with $Q_x(z),P_x(z)$ for $z_0<z\le x+n$.
Since $P_x(z)\ge Q_x(z)$ for every $z$, the minimum of this collection equals
$\min\{Q_x(z):z_0<z\le x+n\}$ when the window extends beyond $z_0$, and equals $Q_x(z_0)+1$
otherwise; in the latter case the stated condition is vacuously true and the letter is indeed
unmatched. The claim follows.
\end{proof}

\begin{lem}[$b_x$ is the last global minimum of the profile]\label{lem:bx}
Let $x\in\XX(S,T)$, let $\mu_x=\min\{Q_x(z):x<z\le x+n\}$, and let $z^{*}$ be the largest $z$
in the window with $Q_x(z)=\mu_x$. Then $L_1^{(x)}(uv)\neq\emptyset$ if and only if
$z^{*}\in S$, and in that case $b_x=z^{*}$.
\end{lem}

\begin{proof}
Suppose $z^{*}\in S$. For every $z>z^{*}$ in the window, $Q_x(z)>\mu_x=Q_x(z^{*})$ by
maximality of $z^{*}$, so by Lemma~\ref{lem:unmatched} the letter at $z^{*}$ is unmatched;
hence $L_1\ne\emptyset$. If $z_0\in S$ is unmatched and $z_0<z^{*}$ then
Lemma~\ref{lem:unmatched} gives $Q_x(z^{*})>Q_x(z_0)\ge\mu_x=Q_x(z^{*})$, absurd; so $z^{*}$ is
the $\prec_x$-smallest unmatched letter, i.e.\ $b_x=z^{*}$.

Suppose instead $z^{*}\notin S$, and assume for contradiction that some $z_0\in S$ is
unmatched. By the displayed argument $z_0\ge z^{*}$, and $z_0\ne z^{*}$, so $z_0>z^{*}$. Let
$z_1$ be the smallest element of $S$ in the window with $z_1>z^{*}$; then $z_1\le z_0$ and every
slot $z$ with $z^{*}<z<z_1$ lies outside $S$, so $P_x(z)=Q_x(z)$ there, and
$P_x(z^{*})=Q_x(z^{*})=\mu_x$ as $z^{*}\notin S$. Since each slot can only decrease the height
before its own low point, $Q_x(z)\le\mu_x$ for all $z^{*}<z\le z_1$; minimality of $\mu_x$
forces equality, and in particular $Q_x(z_1)=\mu_x$ with $z_1>z^{*}$ --- contradicting the
maximality of $z^{*}$. Hence no unmatched letter exists and $L_1=\emptyset$.
\end{proof}

\section{Every crystal edge is an exchange move}

We use one lemma from \cite{Clio2}, namely Step 1 of the proof of \cite[Lem.~4.2]{Clio2}. It is
the only place where length-additivity enters, and we restate it with its proof so that the
present note is self-contained.

\begin{lem}[{\cite[Lem.~4.2, Step 1]{Clio2}}]\label{lem:step1}
Let $(S,T)$ be additive and let $[m,M]$ be a run of $S$ with $m\in T$. Then
$[m,M+1]\subseteq T$.
\end{lem}

\begin{proof}
Let $[a,b]$ be the run of $T$ containing $m$; it is finite because $T$ is proper. Apply the
length-additivity criterion \cite[Lem.~2.3]{Clio2} with $w=u_S$ to the run $[a,b]$ of $T$,
legitimate since $\ell(u_Su_T)=\ell(u_S)+|T|$: it gives $u_S(b+1)>u_S(j)$ for all $j\in[a,b]$,
in particular for $j=m$. Now $m$ is the bottom of the $S$-block $[m,M+1]$, so $u_S(m)=M+1$.
Suppose $b\le M$. If $b=m$ then $b+1\in[m+1,M+1]$, so $u_S(b+1)=m$, and $m>M+1$ is false. If
$m<b\le M$ then $b+1\in[m+1,M+1]$, so $u_S(b+1)=b\le M<M+1=u_S(m)$, again false. Hence
$b\ge M+1$, and since $a\le m$ we get $[m,M+1]\subseteq[a,b]\subseteq T$.
\end{proof}

\begin{thm}\label{thm:contain}
Let $(S,T)$ be additive and let $x\in\XX(S,T)$ with $\et^{(x)}(u_Su_T)\ne\mathbf 0$. Put
$b=b_x$ and let $[m,M]$ be the run of $S$ containing $b$. Then
\begin{enumerate}
\item[(i)] $b+1\notin T$;
\item[(ii)] $[m+1,b]\subseteq T$;
\item[(iii)] $m\notin T$.
\end{enumerate}
Consequently $[m,M]$ is a usable run, $b=e_{[m,M]}$, and
\[
  \et^{(x)}(u_Su_T)\;=\;u_{S\setminus\{b\}}\,u_{T\cup\{m\}}
\]
is exactly the Theorem 4.1 exchange move on the run $[m,M]$. In particular
$\BB(S,T)\subseteq\EE(S,T)$ and every $\et$-image is a Theorem 4.1 move.
\end{thm}

\begin{proof}
Throughout identify $b$ and $m$ with their lifts in the window $(x,x+n]$; this is legitimate
because $[m,b]\subseteq S$ and $x\notin S$, so the whole interval $[m,b]$ avoids the lift of
$x$ and lies in a single window, with $m>x$.

(i) First, $b<x+n$: indeed $b\equiv x+n$ would give $b\equiv x$, impossible since $b\in S$ and
$x\notin S$. So the slot $b+1$ lies in the window. If $b+1\in T$ then slot $b+1$ opens with a
$\texttt{)}$, hence $Q_x(b+1)=P_x(b)-1=Q_x(b)$, and Lemma~\ref{lem:unmatched} shows the letter
at $b$ is matched --- contradicting $b\in L_1^{(x)}$.

(ii) Let $m<j\le b$. Then $j-1\in S$, so $P_x(j-1)=Q_x(j-1)+1$ and hence
$Q_x(j)=Q_x(j-1)+1-[j\in T]$. Thus $Q_x$ is non-decreasing on $[m,b]$ and strictly increases
exactly at those $j$ with $j\notin T$. Suppose some $j\in(m,b]$ has $j\notin T$. Then
$Q_x(j-1)<Q_x(j)\le Q_x(b)$, and $j-1$ lies in the window; but $b=b_x$ satisfies
$Q_x(b)=\mu_x$ by Lemma~\ref{lem:bx}, so $Q_x(j-1)<\mu_x$ contradicts the definition of
$\mu_x$. Hence every $j\in(m,b]$ lies in $T$.

(iii) Suppose $m\in T$. By Lemma~\ref{lem:step1}, $[m,M+1]\subseteq T$. Since $m\le b\le M$ we
get $b+1\in[m+1,M+1]\subseteq T$, contradicting (i).

By (iii) and (i) the run $[m,M]$ satisfies $m\notin T$ and has an index, namely $b$, with
$b+1\notin T$, so it is usable; by (ii) no $j\in[m,b-1]$ has $j+1\notin T$, so $b$ is the
\emph{first} such index, i.e.\ $b=e_{[m,M]}$. Finally $m\notin T$ and $x\notin T$ with $m\ne x$
give $|T|\le n-2$, so \cite[Lem.~4.3]{Clio2} applies to the run $[m,M]$ and its conclusion,
$u_{S\setminus\{b\}}u_{T\cup\{m\}}=u_Su_T$ with the factorisation length-additive, is the
Theorem 4.1 move. Morse--Schilling's operator removes $b$ from $S$ and adds $b-t$, the bottom
of the run of $S$ containing $b$, which is $m$; the two moves coincide.
\end{proof}

\begin{rem}\label{rem:minmatters}
The minimum in $b=\min_{\prec_x}L_1$ is load-bearing, not a canonicalisation of an otherwise
valid family. Part (ii) of the proof is the only step that uses minimality, and it fails for a
general unmatched letter: for $n=3$, $S=\{0,1\}$, $T=\emptyset$, $x=2$ the word is
$\texttt{((}$ and $L_1=\{0,1\}$, but only $b=0$ gives a factorisation --- removing $1$ and
adding $m=0$ yields $u_{\{0\}}u_{\{0\}}=s_0s_0=\mathrm{id}$, of length $0$, not $2$. Over
$3\le n\le7$ the union $\bigcup_x L_1^{(x)}$ fails to be contained in $\EE$ for $2752$ of the
$9093$ additive pairs with $\XX\ne\emptyset$.
\end{rem}

\section{The containment is strict, and where the excess lives}

\begin{thm}\label{thm:strict}
The containment of Theorem~\ref{thm:contain} is strict. The smallest witnesses occur at $n=5$;
there are no witnesses for $n\le4$.
\end{thm}

\begin{proof}
Take $n=5$, $S=\{0,2\}$, $T=\{4\}$, so $u_S=s_0s_2$, $u_T=s_4$ and $v=s_0s_2s_4$ has length
$3=|S|+|T|$: the pair is additive. Note $|S|=2>1=|T|$, so the hypothesis of
\cite[Thm.~4.1]{Clio2} holds and $\et^{(x)}$ fires at every admissible $x$.

\emph{The exchange moves.} $S$ has the two runs $[0,0]$ and $[2,2]$. For $[0,0]$: $0\notin T$
and $1\notin T$, so the run is usable with $e=0$, giving $u_{\{2\}}u_{\{0,4\}}=s_2\cdot s_0s_4
=s_0s_2s_4=v$. For $[2,2]$: $2\notin T$ and $3\notin T$, so the run is usable with $e=2$,
giving $u_{\{0\}}u_{\{2,4\}}=s_0\cdot s_2s_4=v$. Both are length-additive. Hence
$\EE(S,T)=\{0,2\}$.

\emph{The crystal edges.} $\XX(S,T)=\{1,3\}$. For $x=1$ the window is $\{2,3,4,5\}\cup\{6\}$,
whose word is $\texttt{(}\,\texttt{)}\,\texttt{(}$ at slots $2,4,5$; the letters at $2$ and $4$
match and the letter at $5\equiv0$ is unmatched, so $L_1=\{0\}$ and $b_1=0$. For $x=3$ the
window is $\{4,\dots,8\}$, whose word is $\texttt{)}\,\texttt{(}\,\texttt{(}$ at slots
$4,5,7$; here $L_1=\{0,2\}$, and since \eqref{eq:order} for $x=3$ reads
$2\succ1\succ0\succ4$, the $\prec_3$-minimum is $b_3=0$. Hence $\BB(S,T)=\{0\}\subsetneq
\{0,2\}=\EE(S,T)$: the exchange move on the run $[2,2]$ is an edge of no Morse--Schilling
crystal.

The mechanism is visible in the profile. Normalising $P(-1)=0$ one finds $Q(0)=0$, $Q(1)=1$,
$Q(2)=1$, $Q(3)=2$, $Q(4)=1$ and $d=|S|-|T|=1$, so $Q(5)=Q(0)+1=1=Q(2)$. In \emph{every}
window of length $n$ containing both, the slots $2$ and $5$ carry equal profile values and
slot $5\equiv0$ comes later; by Lemma~\ref{lem:bx} the operator selects the \emph{last}
minimiser, so slot $2$ is a minimiser that is never the last one. Exhaustive verification for
$n\le4$ (\S\ref{sec:verify}) shows $\BB=\EE$ there.
\end{proof}

\begin{thm}[Local versus global minimality]\label{thm:charac}
Let $(S,T)$ be additive with $\XX(S,T)\ne\emptyset$ and let $c\in S$. Then:
\begin{enumerate}
\item $c\in\EE(S,T)$ if and only if, with $m$ the bottom of the run of $S$ containing $c$, one
has $m\notin T$, $[m+1,c]\subseteq T$ and $c+1\notin T$. Equivalently: $Q$ is constant on
$[m,c]$, $Q$ increases strictly at $c+1$, and $m\notin T$. This is a condition on the single
run of $c$ --- it is \emph{local}.
\item $c\in\BB(S,T)$ if and only if $c\in\EE(S,T)$ and there exists $x\in\XX(S,T)$, lifted to
$\hat x$ with $c-n\le\hat x<c$, such that
\[
  Q(z)\ \ge\ Q(c)\quad (\hat x<z<c)
  \qquad\text{and}\qquad
  Q(z)\ >\ Q(c)\quad (c<z\le \hat x+n).
\]
Equivalently, writing
\[
  F(c)=\min\{z>c: Q(z)\le Q(c)\},\qquad G(c)=\max\{z<c: Q(z)<Q(c)\}
\]
(with $F=+\infty$, $G=-\infty$ if no such $z$ exists), $c\in\BB(S,T)$ if and only if the
integer interval $[\,\max(G(c),c-n),\ \min(F(c)-n-1,\,c-1)\,]$ contains a lift of an element of
$\XX(S,T)$. This is a condition on a full period --- it is \emph{global}.
\end{enumerate}
\end{thm}

\begin{proof}
(1) is Definition~\ref{def:E} unwound: $[m,M]$ is usable with $e_{[m,M]}=c$ exactly when
$m\notin T$, $j+1\in T$ for $m\le j<c$ and $c+1\notin T$; the first of these two conditions is
$[m+1,c]\subseteq T$. The reformulation in terms of $Q$ is the computation already made in the
proof of Theorem~\ref{thm:contain}(ii): for $m<j\le c$ one has $Q(j)=Q(j-1)+1-[j\in T]$, so
$[m+1,c]\subseteq T$ says precisely that $Q$ is constant on $[m,c]$, and $c+1\notin T$ says
$Q(c+1)=P(c)=Q(c)+1$.

(2) By Lemma~\ref{lem:bx}, for $x\in\XX(S,T)$ we have $b_x=c$ if and only if $c$ lies in the
window $(x,x+n]$ and is the largest slot of that window at which $Q_x$ attains its minimum over
the window. Since $Q_x$ and $Q$ differ by a constant on the window, this says exactly
$Q(z)\ge Q(c)$ for $\hat x<z<c$ and $Q(z)>Q(c)$ for $c<z\le\hat x+n$, which is the displayed
condition. The first clause holds for all $z\in(\hat x,c)$ iff $\hat x\ge G(c)$, and the second
holds for all $z\in(c,\hat x+n]$ iff $\hat x+n<F(c)$, i.e.\ $\hat x\le F(c)-n-1$; together with
$c-n\le\hat x<c$ this is the stated interval. Finally, $b_x=c$ forces $c\in\EE(S,T)$ by
Theorem~\ref{thm:contain}.
\end{proof}

\begin{rem}
This is what the pairing buys. A crystal operator must be a \emph{function}: it must select one
of the available moves, and it must do so in a way compatible with $\tilde f_1$, with the
Stembridge axioms, and with the weight. Morse--Schilling's selection rule is ``take the last
global minimum of the profile over the window'', and globality is precisely what the
$x$-pairing computes --- a pairing is a statement about the whole word, a run is not. Theorem
4.1 asks only for local minimality and therefore sees moves that no global rule can reach: in
the example of Theorem~\ref{thm:strict}, slot $2$ ties with a later slot in every window, and
a rule that must break ties consistently can never choose it. Conversely, having no parameter,
the exchange move makes no claim to canonicity: $\EE(S,T)$ is a set, not an operator, and
Theorem 4.1 asserts only that it is nonempty when $|S|>|T|$, which is all
\cite[Cor.~4.5]{Clio2} consumes.
\end{rem}

\section{The case $\XX(S,T)=\emptyset$}\label{sec:emptyX}

When $S\cup T=\Z/n$ there is no admissible $x$ and Definition~\ref{def:word} does not apply.
Morse--Schilling handle this in \cite[\S\texttt{two factor case}]{MS}: they bracket $i\in S$ with $i+1\in T$, show
(their \texttt{lemma.existence of i}) that some index is in ``Case (3)'', i.e.\ lies in a maximal block
\[
  \{b-1,\dots,b-t\}\subseteq S,\ b\notin S,\qquad \{b-t+1,\dots,b\}\subseteq T,\ b-t\notin T,
\]
and then remove the bracketed pair $(b-1,b)$ and proceed with $x=b$.

Two points of honesty. First, the prescription ``proceed with $x=b$'' leaves it unstated
whether the offset $t$ of \cite[Def.~\texttt{crystal operators}(i)]{MS} is computed against $\cont(u)$ or against the
reduced content $\cont(u)\setminus\{b-1\}$; we implemented both readings. Second, the choice of
Case-(3) index $b$ is not in general unique, so this construction, like the main one, defines a
\emph{family} of operators, and we take the union over the family. With those two conventions
fixed, exhaustive computation for $3\le n\le7$ gives, over all $1636$ additive pairs with
$\XX=\emptyset$:
\begin{itemize}
\item both readings produce the same set of moves in every case;
\item every move so produced is a genuine length-additive factorisation of $u_Su_T$ ($1636$ of
$1636$ pairs, checked by comparing window notations directly);
\item every move so produced lies in $\EE^{\mathrm{mv}}(S,T)$ ($1636$ of $1636$), and the
containment is again strict, failing to be an equality in $7$ pairs, all at $n=7$.
\end{itemize}
So Theorem~\ref{thm:contain} and Theorem~\ref{thm:strict} are not artefacts of restricting to
$\XX\ne\emptyset$: the conclusion is uniform over all two-factor affine factorisations. We
record the $\XX=\emptyset$ statement as computational, not proved --- see \S\ref{sec:gaps}.

\section{Computational verification}\label{sec:verify}

All code is in \texttt{proofs/code-q220/} of the commit named on the title page. Affine
permutations are held in window notation and lengths computed by Shi's inversion formula, not
by the block model, so the block model of \cite[Lem.~2.2]{Clio2} is tested rather than assumed.

\paragraph{Instrument validation.} The implementation of \cite[\S\texttt{the crystal operators}]{MS} reproduces all
three pairings printed in \cite[Ex.~\texttt{ex:pairing}]{MS} exactly, including the pair labels
($n=14$, $x=10$: pairs $(5,6),(2,4),(12,13)$, $L_1=\{8,9\}$; $n=14$, $x=3$: pairs
$(12,13),(9,11),(8,0),(5,6)$; $n=5$, $x=2$: pair $(0,1)$), and the implementation of
\cite[Def.~\texttt{crystal operators}(i)]{MS} reproduces all three $\et$-images printed in \cite[the example following Def.~\texttt{crystal operators}]{MS}
exactly. No parameter was fitted.

\paragraph{The lemmas.} Lemma~\ref{lem:bracket} was checked on every admissible cut of every
additive pair for $3\le n\le7$ ($17{,}368$ cuts, $0$ disagreements between the Morse--Schilling
procedure and bracket matching). Lemma~\ref{lem:bx} was checked on the same $17{,}368$ cuts,
including the ``only if'' half ($0$ failures). Lemma~\ref{lem:step1} was checked on all $3{,}742$
instances of its hypothesis ($0$ failures).

\paragraph{Theorem~\ref{thm:contain}.} Clauses (i), (ii), (iii), the identification
$b=e_{[m,M]}$, and the equality of the two moves were each checked separately at all
$10{,}877$ firings of $\et^{(x)}$ over $3\le n\le7$: $0$ failures in each of the five tests.
Independently, all $9{,}120$ Theorem 4.1 moves were re-checked to be genuine length-additive
factorisations with proper supports ($0$ failures); this tests the implementation of
\cite[Thm.~4.1]{Clio2}, not the theorem.

\paragraph{Theorem~\ref{thm:charac}.} The criterion of part (2) was evaluated independently of
the pairing, from the profile alone, and compared with $\BB(S,T)$ computed from
\cite[Def.~\texttt{crystal operators}]{MS}: agreement on $9{,}093$ of $9{,}093$ additive pairs with $\XX\ne\emptyset$.

\paragraph{The counts.} For $3\le n\le7$:
\begin{center}
\begin{tabular}{rrrrrrr}
\hline
$n$ & additive pairs & $\XX=\emptyset$ & $|\EE|$ & $|\BB|$ & pairs with $\BB=\EE$ & excess letters \\
\hline
3 & 34   & 9    & 15   & 15   & 34   & 0 \\
4 & 143  & 34   & 80   & 80   & 143  & 0 \\
5 & 561  & 115  & 375  & 360  & 546  & 15 \\
6 & 2122 & 365  & 1650 & 1518 & 1990 & 132 \\
7 & 7869 & 1113 & 7000 & 6146 & 7015 & 854 \\
\hline
\end{tabular}
\end{center}
Here $|\EE|$ and $|\BB|$ are summed over all additive pairs, $\BB$ being computed with the
$\XX=\emptyset$ convention of \S\ref{sec:emptyX} where needed, and ``excess letters'' counts
$\sum|\EE\setminus\BB|$. $\BB\subseteq\EE$ held in all $10{,}729$ pairs.

\begin{rem}
No coverage \emph{fraction} is quoted anywhere in this note. Two independent implementations of
the Morse--Schilling pairing agreed on the forward direction on 21 September and disagreed on a
converse coverage fraction; that discrepancy is unresolved, so only integer counts, each with
its defining predicate spelled out, appear above.
\end{rem}

\section{What this consumes, and what remains open}\label{sec:gaps}

\paragraph{Consumed.} \S\texttt{the crystal operators}, Def.~\texttt{crystal operators} and \S\texttt{two factor case} of \cite{MS}, read at source; Lemmas
2.2, 2.3, 4.2 (Step 1) and 4.3 and Theorem 4.1 of \cite{Clio2}. Nothing is consumed from
\cite{Lam2008}, from the theory of dual $k$-Schur functions, or from
\cite[Thm.~\texttt{theorem.crystal}]{MS} (that $B(w)$ is a crystal): Theorem~\ref{thm:contain} is a statement about
the operator $\et$, not about the crystal axioms it satisfies.

\paragraph{Gaps.}
\begin{enumerate}
\item[(G1)] \S\ref{sec:emptyX} is \emph{computational, not proved}. The obstruction is precise:
Theorem~\ref{thm:contain}(iii) invokes Lemma~\ref{lem:step1}, which needs $(S,T)$ additive,
whereas the $\XX=\emptyset$ construction pairs the \emph{reduced} pair
$(S\setminus\{b-1\},T\setminus\{b\})$, which need not be additive. I do not know whether
Theorem~\ref{thm:contain} holds there for every choice of Case-(3) index; it holds in all
$1636$ instances with $n\le7$.
\item[(G2)] The reading of ``$t$ maximal'' and of ``proceed with $x=b$'' in \cite[\S\texttt{two factor case}]{MS} is
mine, not theirs. Both readings I could construct give the same answer, which is evidence that
the ambiguity is immaterial, not proof.
\item[(G3)] Theorem~\ref{thm:charac}(2) characterises $\BB$ by a condition that still quantifies
over $\XX$. A characterisation intrinsic to $(S,T)$ --- one that predicts $|\EE\setminus\BB|$
from the run structure without sliding a window --- would be better, and I do not have one. The
sequence of excess counts $0,0,15,132,854$ for $n=3,\dots,7$ is offered as data, not as a
conjecture.
\item[(G4)] Theorem~\ref{thm:contain} says every $\et$-edge is an exchange move; it says nothing
about $\tilde f_1$ or $\tilde s_1$. Since $\tilde f_1$ is inverse to $\et$ where both are
defined, one expects the corresponding statement for $\tilde f_1$ to follow from the dual
exchange move obtained in the proof of \cite[Cor.~4.5]{Clio2} by conjugating with the
length-preserving anti-automorphism $\Psi$; I have not checked this and do not claim it.
\end{enumerate}

\begin{thebibliography}{9}
\bibitem{Clio2} Clio, \emph{An exchange move for cyclically decreasing factorisations, and the
reduction of WZZ Problem 5.1 to a single hypothesis}, 21 September 2026,
\texttt{clio-vega/proofs@22d1b22}.
\bibitem{Lam2006} T.~Lam, \emph{Affine Stanley symmetric functions}, Amer.\ J.\ Math.\ 128
(2006), 1553--1586.
\bibitem{Lam2008} T.~Lam, \emph{Schubert polynomials for the affine Grassmannian}, J.\ Amer.\
Math.\ Soc.\ 21 (2008), 259--281.
\bibitem{MS} J.~Morse and A.~Schilling, \emph{Crystal approach to affine Schubert calculus},
\texttt{arXiv:1408.0320}.
\end{thebibliography}

\end{document}
